\documentclass[10pt,a4paper]{amsart}
\usepackage[margin=19mm]{geometry}
\usepackage{amsmath,amssymb,mathtools}
\usepackage[scr=rsfs,cal=euler]{mathalfa}
\usepackage{xfrac}
\usepackage[unicode,hidelinks,bookmarks=false]{hyperref}
\newcommand{\ie}{i.e.\ }
\newcommand{\cf}{cf.\ }
\newcommand{\resp}{respectively\ }
\newcommand{\Subsection}{\subsection}
\providecommand{\nobreakdash}{\nobreak\hbox{-}}
\setlength{\emergencystretch}{2em}
\multlinegap0pt
\usepackage{bm}
\usepackage{bbm}
\usepackage{dsfont}
\usepackage{xspace}
\usepackage{cancel}
\usepackage{mathtools}
\usepackage{amsthm}
\makeatletter
\@ifundefined{prop}{\newtheorem{prop}{Proposition}}{}
\@ifundefined{lem}{\newtheorem{lem}[prop]{Lemma}}{}
\makeatother
\newcommand{\CC}{\mathbb{C}}
\newcommand{\C}{\mathcal{C}}
\newcommand{\RR}{\mathbb{R}}
\newcommand{\ZZ}{\mathbb{Z}}
\newcommand{\m}{\mathfrak{m}}
\newcommand{\n}{\mathfrak{n}}
\title[Entire Nodal Solutions with Prescribed Symmetry to Caffarell]{Entire Nodal Solutions with Prescribed Symmetry to Caffarelli-Kohn-Nirenberg Equations}
\author{Edward Chernysh}
\date{Source: arXiv:2511.11543v1 (CC BY 4.0)}
\begin{document}
\maketitle

\noindent\emph{Source and licence.} Entire Nodal Solutions with Prescribed Symmetry to Caffarelli-Kohn-Nirenberg Equations. Version arXiv:2511.11543v1 (CC BY 4.0), distributed under the Creative Commons Attribution 4.0 International licence, as recorded in the arXiv licence metadata for this exact version and shown on the abstract page https://arxiv.org/abs/2511.11543v1. Full rights notice: licenses/arxiv-2511.11543-CC-BY-4.0.txt.\par\smallskip

\noindent\textit{Editorial scope.} Counted unit: the Prop whose source label is \texttt{prop:distinct2} in the article (source0.tex lines 498--500, source label \texttt{prop:distinct2}), delivered together with the article's own complete original proof of it (source0.tex lines 501--564, one \texttt{proof} environment of 64 lines). No mathematics is written, repaired or re-derived: statement and proof are the article's own text line for line. Required context retained uncounted: eq:primeCondition (lines 95--98); lem:isCode (lines 411--417). Every definition, display and label of those lines is retained unchanged. Omitted from this package and not redistributed: 1--94, 99--410, 418--497, 565--1096. Nothing inside a retained excerpt is omitted or altered: every retained body is delivered exactly as the article's lines are written, so no omission receipt of any kind accompanies this package. Citations: the retained excerpts carry no citation command at all, so no bibliography entry of the article is transcribed and no reference is redistributed. Reference adaptation: none was needed and none was made; every reference inside the delivered unit resolves inside it, so no unresolved reference or citation is produced. Printed slips kept literal: no printed word, symbol, hypothesis or number of the counted statement or of its proof was altered. Layout and class: the article's own class is \texttt{amsart} with packages including amsmath, amssymb, amsthm, mathtools, enumitem, xcolor, hyperref; the wrapper is the lane's pinned \texttt{amsart} editorial wrapper at 10pt on A4, which restates the theorem-like environments the retained text needs and adds the article's own macro definitions that the retained text needs (\texttt{\textbackslash C}, \texttt{\textbackslash CC}, \texttt{\textbackslash RR}, \texttt{\textbackslash ZZ}, \texttt{\textbackslash m}, \texttt{\textbackslash n}), with extra wrapper packages (bm, bbm, dsfont, xspace, cancel, mathtools). The delivered numbering is the unit's own. Assets: no retained span contains a figure environment, a graphics-inclusion command, a PDF-page inclusion, a table or a TikZ picture; no image file is copied into this package, the package declares no asset dependency and no image file is redistributed with it. Licence: version arXiv:2511.11543v1 of this article is distributed under the Creative Commons Attribution 4.0 International licence, as recorded in the arXiv licence metadata for this exact version and shown on the abstract page https://arxiv.org/abs/2511.11543v1. A full rights notice is delivered at licenses/arxiv-2511.11543-CC-BY-4.0.txt. Characters: every retained excerpt is pure ASCII, so the delivered engine, XeLaTeX with its default Latin Modern text font, reports no missing character.
	\begin{align}
		&0 < 2\chi_\alpha  + \textstyle\sum_{j=1}^k m_j(j+1) \le n/2,\label{eq:primeCondition}\\
		&\textstyle 2\chi_\alpha + \sum_{j=1}^k m_j(j+1) \ne (n-1)/2.\label{eq:primeCondition2}
	\end{align}

	\begin{lem}\label{lem:isCode}
		Let $n \ge 4$ and $(\alpha,\m)$ be a pair for which \eqref{eq:primeCondition} holds. Suppose $1 \le j \le k$ is an index such that $m_j >0$ and fix $\ell \in \{1,\dots, m_j\}$. Assume $u$ is $(\alpha,\m)$-symmetric and let $\C$ be the set of all binary codewords $c \in \ZZ_2^{j+1}$ such that
		\[
		u(R_c^{\alpha,\m, j,\ell}(\theta)x) = u(x)
		\]
		for all $x \in \RR^n$ and $\theta \in [0,2\pi)$. Then, $\C$ is a $(j+1)$-code.
	\end{lem}

	\begin{prop}\label{prop:distinct2}
		Let $n \ge 4$ and suppose $(\alpha, \m)$ and $(\beta, \n)$ are two tuples, each satisfying \eqref{eq:primeCondition}. Let $u : \RR^n \to \RR$ be $(\alpha,\m)$-symmetric and $(\beta,\n)$-symmetric. If $\alpha \ne \beta$, then $u$ vanishes identically.
	\end{prop}

	\begin{proof}
		Since $\alpha \ne \beta$, we may assume without loss of generality that $\alpha < \beta$. In this case, we must distinguish between $\alpha = 0$ and $\alpha > 0$. In the latter case, we have 
		\(
		\varrho_\alpha = \varrho_{\beta}^{2^{\beta-\alpha}}
		\)
		implying that, for all $x \in \RR^n$, 
		\begin{align*}
			u(x) = u(\varrho_{\beta}^{2^{\beta-\alpha}}(x))
			= u(\varrho_\alpha(x))
			&= -u(x),
		\end{align*}
		where the first equality uses the $(\beta,\n)$-symmetry of $u$ and the final equality invokes the $(\alpha,\m)$-symmetry. Consequently, we infer that $u \equiv 0$.
		
		If instead $\alpha = 0$, then observing that $\varrho_{\beta}^{2^{\beta}} = \rho_1$ yields a similar contradiction provided $m_1 \ne 0$. Indeed,
		\begin{align*}
			u(x) = u(\varrho_{\beta}^{2^{\beta}}(x))
			= u(\rho_1(x))
			&= -u(x)
		\end{align*}
		whence $u \equiv 0$.
		It remains to treat the setting where $0 = \alpha < \beta$ and $m_1 = 0$. In this case,  let $j > 1$ be minimal such that $m_j \ne 0$. Denote by $(z_1,z_2,\dots, z_{j+1})$ the complex coordinates acted upon by the first copy of $\Gamma_j$ in $G_{\alpha, \m}$. Then, for any $x = (z_1,\dots, z_{j+1}, w) \in \RR^n \cong \CC^{j+1} \times \RR^{n-2(j+1)}$, there holds
		\begin{equation}\label{eq:rotateAll}
			u(x) = u(e^{i\theta}z_1, \dots, e^{i\theta} z_{j+1}, w)
		\end{equation}
		for any $\theta \in \RR$, by virtue of the $(\alpha,\m)$-symmetry. Owing to the $(\beta, \n)$-symmetry, we also have
		\begin{align*}
			u(x) = u(e^{i\theta}z_1, e^{-i\theta} z_{2}, z_3, \dots, z_{j+1}, w)
		\end{align*}
		for all angles $\theta \in \RR$. By following the argument in Lemma \ref{lem:isCode} (i.e. conjugating the above rotations by the copy of $\rho_j$ in $G_{\alpha,\m}$ acting on the coordinates $(z_1,\dots,z_{j+1})$), we see that 
		\begin{align}
			u(x) &= u(e^{i\theta}z_1, e^{-i\theta} z_{2}, z_3, \dots, z_{j}, z_{j+1}, w) \tag{E.1}\\
			&= u(z_1, e^{i\theta} z_{2}, e^{-i\theta}z_3, z_4, \dots, z_{j}, z_{j+1}, w)\tag{E.2}\\
			&= u(z_1, z_2, e^{i\theta} z_{3}, e^{-i\theta}z_4, \dots, z_{j}, z_{j+1}, w)\tag{E.3}\\
			&\,\,\,\vdots\tag*{}\\
			&= u(z_1, z_2, z_{3}, z_4, \dots, e^{i\theta}z_{j}, e^{-i\theta}z_{j+1}, w) \tag{E.j}
		\end{align}
		Now, given $\theta$, we sequentially apply (E.1) with $\theta/(j+1)$, (E.2) with $2\theta/(j+1)$, (E.3)  with $3\theta/(j+1)$, and so forth, until we apply (E.j) with $j\theta/(j+1)$. Each such rotation preserves the sign of $u$, whence
		\[
		u(x) = u(e^{i\theta/(j+1)}z_1, e^{i\theta/(j+1)} z_2, \dots, e^{i\theta/(j+1)} z_{j}, e^{-ij\theta/(j+1)}z_{j+1}, w).
		\]
		Then, applying \eqref{eq:rotateAll} with $-\theta/(j+1)$, we obtain
		\[
		u(x) = u(z_1, z_2, \dots, z_{j}, e^{-i\theta}z_{j+1}, w).
		\]
		Consequently, $u$ is rotation invariant in the complex coordinate $z_{j+1}$. Following once more the conjugation argument from Lemma \ref{lem:isCode}, it follows that $u$ is rotation invariant in each complex coordinate of $(z_1,\dots, z_{j+1})$. 
		
		Next, by using that $u$ is $(\beta,\n)$-symmetric with $\beta > 0$, we obtain for every point $x  = (z,w)\in \RR^n \cong \CC^{j+1} \times \RR^{n-2(j+1)}$ that 
		\begin{align*}
			u(x) = u(\varrho_{\beta}^{2^\beta}(x)) &= u(-\bar{z_2}, \bar{z_1}, z_3, z_4, \dots, z_{j}, z_{j+1}, w)\\
			&= u({z_2}, {z_1}, z_3, z_4, \dots, z_{j}, z_{j+1}, w),
		\end{align*}
		where we have used that $u$ is invariant under rotation in each coordinate of $z$ for this last step. Again, by conjugation with respect to $\rho_j$, the same argument as before (see also Lemma \ref{lem:isCode}) implies that any two adjacent coordinates of $z = (z_1,\dots,z_{j+1})$ may be swapped while preserving the value of $u(z)$. Iterating this principle yields the identity
		\begin{equation}\label{eq:distinctEvenCase}
			u(x) = u(z_{j+1}, z_1, z_2, \dots, z_j, w).
		\end{equation}
		On the other hand, the $(\alpha, \m)$-symmetry of $u$ together with the rotational invariance in each coordinate of $z$ ensures that
		\begin{align*}
			-u(x) = u(\rho_j (x)) &= u(-\overline{z_{j+1}}, \bar{z_1}, \dots, \bar{z_j},w)\\
			&=u(z_{j+1}, z_1, \dots, z_j, w).
		\end{align*}
		Combining this last equation with \eqref{eq:distinctEvenCase}, we conclude that $u(x) = 0$. 
		
		In all cases we have shown that $u \equiv 0$, which is what had to be shown.
	\end{proof}

\end{document}
